% from main.tex
%%%%%%%%%%  scan idx 700  |  folio 687  %%%%%%%%%%
isolées sont des invariants géométriques tout à fait non triviaux, et qui auront sans doute leur rôle à jouer dans l'étude ``arithmétiques'' de ces anneaux locaux.

Bien entendu, une fois en possession d'une notion raisonnable d'anneau de Chow pour un schéma noethérien régulier admettant un Module inversible ample, il y a lieu de généraliser à ce cas les relations établies resp. conjecturées au no. 4 entre cet anneau et l'anneau $\Gtop^\bullet(X)$, et à formuler en termes de cet anneau la question de la ``formule de Riemann--Roch sans dénominateurs'' énoncée au no. 3.

\subsection*{Bibliographie de l'Exposé XIV}
\begin{enumerate}[label={[\arabic*]}]
\item M. F. Atiyah et F. Hirzebruch, \emph{The Riemann--Roch theorem for analytic embeddings}, Topology, vol. 1, pp. 151--166 (1962).
\item A. Borel et J.-P. Serre, \emph{Le théorème de Riemann--Roch}, Bull. Soc. Math. France, vol. 86, pp. 97--136 (1958).
\item R. Bott et H. Samelson, Amer. Journ. of Math. vol. 80, p. 1004 (1958).
\item C. Chevalley, \emph{Les classes d'équivalence rationelle}, I et II, Séminaire à l'ENS 1958, exposés 2 et 3.
\item C. Chevalley, \emph{Classification des groupes de Lie algébriques}, Séminaire à l'ENS 1956/58.
\item A. Dold et D. Puppe, \emph{Homologie nicht-additiver Funktoren. Anwendungen}, Annales de l'Inst. Fourier, t. XI, pp. 201--312 (1961).
\item H. Grauert, \emph{Ein Theorem der analytischen Garbentheorie}, Pub. Math. no. 5, pp. 5--64 (1960).
\item A. Grothendieck, \emph{Classes de faisceaux et théorème de Riemann--Roch}, notes multigraphiées par les soins de l'Institute for Advanced Study, 1957 (travail cité [RRR] dans le présent exposé, et qui est reproduit dans notre séminaire sous forme d'Appendice au Chap. O).
\item A. Grothendieck, \emph{La théorie des classes de Chern}, Bull. Soc. Math. France, vol. 86, pp. 137--154 (1958).
\item A.
